\chapter{The Biquaternionic Bulk and Relativistic Kinematics}

The construction of a holographic spacetime requires a native coordinate system capable of unifying 4D Minkowski geometry with spinorial degrees of freedom. In this chapter, we establish the biquaternion algebra as the primary coordinate engine of the bulk and demonstrate its formal equivalence to the Lorentz group $SL(2, \mathbb{C})$ and special relativity.

\section{The Biquaternion Algebra $\mathbb{B}$}
The biquaternion algebra, denoted $\mathbb{B} \cong \mathbb{C} \otimes \mathbb{H}$, is the complexification of Hamilton's quaternions. It is isomorphic to the algebra of $2 \times 2$ complex matrices, $M_2(\mathbb{C})$. A general element $X \in \mathbb{B}$ is represented as:
\begin{equation}
    X = a_0 \mathbf{1} + a_1 \mathbf{i} + a_2 \mathbf{j} + a_3 \mathbf{k}
\end{equation}
where $a_\mu \in \mathbb{C}$. In our dual-witness framework, we utilize the representation of $\mathbb{B}$ via Pauli matrices $\{\sigma_1, \sigma_2, \sigma_3\}$, where $\mathbf{i} \mapsto -i\sigma_1$, $\mathbf{j} \mapsto -i\sigma_2$, and $\mathbf{k} \mapsto -i\sigma_3$.

\section{Minkowski Spacetime as a Hermitian Subspace}
We map a physical event in 3+1 dimensional Minkowski spacetime $(t, x, y, z)$ to a Hermitian biquaternion $X$ by:
\begin{equation}
    X = t I_2 + x \sigma_1 + y \sigma_2 + z \sigma_3 = 
    \begin{pmatrix} 
    t + z & x - iy \\ 
    x + iy & t - z 
    \end{pmatrix}
\end{equation}
The central result of this mapping is the identity between the matrix determinant and the invariant spacetime interval.
\begin{theorem}[Minkowski-Biquaternion Identity]
For a spacetime biquaternion $X$, the determinant evaluates to the Minkowski metric:
\begin{equation}
    \det(X) = t^2 - x^2 - y^2 - z^2
\end{equation}
\end{theorem}
\begin{proof}[Verification]
Formally proven in \texttt{MinkowskiBiquaternion.lean} via \texttt{det\_spacetime\_minkowski} and verified analytically in \texttt{minkowski\_biquaternion\_sympy.py}.
\end{proof}

\section{The Lorentz Group and KAN Decomposition}
Lorentz transformations are realized as elements of the group $SL(2, \mathbb{C})$, the double cover of $SO^+(1,3)$. Any transformation $\Lambda \in SL(2, \mathbb{C})$ acts on the spacetime matrix via conjugation: $X' = \Lambda X \Lambda^\dagger$. 

To understand the holographic depth, we apply the \textit{Iwasawa (KAN) decomposition}, which factors any $\Lambda$ into:
\begin{itemize}
    \item \textbf{K (SU(2))}: Spherical spatial rotations.
    \item \textbf{A ($\mathbb{R}_+^\times$)}: Radial scaling boosts (Holographic depth).
    \item \textbf{N ($\mathbb{C}$)}: Nilpotent boundary translations (Lightcone coordinates).
\end{itemize}

\section{Modular Hamiltonians and Bisognano-Wichmann Boosts}
Following the Bisognano-Wichmann theorem, the modular flow of the boundary is generated by a Lorentz boost. We define the modular Hamiltonian $K$ as a traceless biquaternion. 
\begin{theorem}[Lorentz Invariance of the Flow]
The transformation $X' = e^{\frac{\eta}{2}\sigma_1} X e^{\frac{\eta}{2}\sigma_1}$ represents a physical boost with rapidity $\eta$, preserving $\det(X') = \det(X)$.
\end{theorem}
The formalization in \texttt{lorentz\_invariance} (Lean 4) demonstrates that this boost is not merely a symmetry of the equations, but a fundamental property of the biquaternion algebra's unit group. This establishes the continuous bulk kinematics as the stable foundation for the subsequent boundary scaling-limit analysis.
